\chapter{Introduction}

\begin{cle}
\textbf{Question centrale.} Comment traduire un niveau de non-conformité à DORA, même
hypothétique ou graduel, en une distribution de pertes puis en une charge de capital ? Le
mémoire ne pose pas la question \emph{l'entité est-elle conforme aujourd'hui}, mais
\emph{combien coûte en capital de ne pas se conformer}.
\end{cle}

Le règlement DORA, entré en application en janvier 2025, impose aux entités financières la
maîtrise de cinq domaines de résilience opérationnelle numérique : la gouvernance du risque
lié aux technologies (P1), la gestion et la notification des incidents (P2), les tests de
résilience (P3), la gestion du risque lié aux prestataires tiers (P4) et le partage
d'informations sur les cybermenaces (P5). Ces cinq piliers ne sont pas des cases
indépendantes à cocher : ce sont des domaines de \emph{contrôle}, et la défaillance de l'un
dégrade la capacité de l'entité à tenir les autres. Une gouvernance défaillante retarde la
détection des incidents ; une dépendance non maîtrisée à un prestataire ouvre une porte que
les tests n'ont pas explorée. La non-conformité ne reste pas dans son pilier, elle se
propage.

% SOURCES-SCRIPTS: 29

\section{Le risque cyber en quelques cas}
% HARNAIS-HORS-SECTION: section descriptive sans valeur chiffree ; les seuls nombres
% sont des millesimes d'incidents publics.

Le risque cyber désigne les pertes liées à une défaillance des systèmes d'information, qu'elle
soit malveillante ou non. Quatre familles d'incidents se distinguent : le rançongiciel, qui
chiffre les systèmes et conditionne leur restitution à une rançon ; la fuite ou le vol de
données ; le déni de service, qui rend un service indisponible ; et la défaillance d'un
prestataire ou d'un logiciel tiers, qui atteint simultanément tous ses clients. DORA ne se limite
pas aux attaques : il vise tout incident lié aux technologies de l'information et de la
communication, y compris une panne sans intention de nuire. Trois cas publics montrent pourquoi
les cinq piliers ne cèdent pas indépendamment.

\paragraph{NotPetya, 2017.} Un logiciel destructeur diffusé par la mise à jour d'un logiciel
comptable ukrainien s'est propagé aux réseaux de groupes internationaux et en a paralysé
l'activité pendant plusieurs semaines. L'incident naît chez un prestataire (P4), se propage faute
de cloisonnement et de gouvernance des accès (P1), et met à l'épreuve la gestion de crise (P2).
Son traitement assurantiel a en outre donné lieu à des contentieux sur l'exclusion de guerre des
contrats.

\paragraph{MOVEit, 2023.} Une vulnérabilité d'un logiciel de transfert de fichiers, exploitée par
un groupe criminel, a permis de dérober des données chez des milliers d'organisations, souvent
atteintes à travers leurs propres prestataires. C'est un choc de cause commune : une seule
défaillance tierce produit de nombreux sinistres simultanés. Le mémoire l'utilise comme
expérience naturelle au chapitre~\ref{ch:identifiabilite}.

\paragraph{CrowdStrike, 2024.} Une mise à jour défectueuse d'un logiciel de sécurité a bloqué des
millions de postes de travail dans le transport aérien, la banque et la santé. Aucune attaque
n'est en cause : l'incident relève de la dépendance à un prestataire (P4) et de l'absence de
tests progressifs avant déploiement (P3), et il entre pleinement dans le champ de DORA.

Ces trois cas ont une structure commune : un incident naît sur un domaine de contrôle et en
dégrade d'autres, dans un ordre. C'est cette structure que le mémoire modélise.

\section{Le vide que ce mémoire comble}
Aucun dispositif prudentiel ne traduit aujourd'hui ce niveau de maîtrise en capital. La
Formule Standard de Solvabilité~II charge le risque opérationnel par un pourcentage de primes
ou de provisions, entièrement aveugle au fait qu'une entité soit exemplaire ou
défaillante sur DORA. Une copule capturerait bien la co-occurrence des pertes, mais
symétriquement : elle ne distingue pas une défaillance qui \emph{entraîne} les autres d'une
défaillance qui les \emph{subit}. Or c'est précisément cette asymétrie, la direction de la
propagation, qui devrait commander la priorité de remédiation d'un assureur. C'est le vide
que le mémoire vient combler, sans prétendre à ce que la donnée ne permet pas.

\section{Le marché français de la cyberassurance}
\label{sec:lucy}

Le vide décrit ci-dessus n'est pas une particularité de la Formule Standard : il se constate sur
le marché lui-même. La seule observation continue du marché français en donne la mesure, et elle
impose au mémoire deux contraintes, l'une sur le choix de la base de calibration, l'autre sur la
façon de présenter un écart de capital.

\paragraph{Le cadre d'observation.} L'étude LUCY de l'AMRAE, dont l'édition $2026$ porte sur
l'exercice $2025$, observe le marché français à partir de $20\,996$ polices et $1\,251$ sinistres
déclarés, recueillis auprès de douze courtiers et d'un assureur \citep{amraeLucy2026}. La série
est engagée depuis $2019$, ce qui donne sept exercices comparables et permet de distinguer un
régime d'un accident d'exercice. L'auteur de ce mémoire a co-écrit, avec son tuteur, l'analyse
actuarielle de cette édition publiée par Nexialog Consulting \citep{RapiorKaddouri2026}. Cette
source est employée ici sous un statut précis, celui de \emph{citation externe non
recalculable} : ses valeurs sont reprises telles qu'elle les publie, sans que le mémoire soit en
mesure de les recalculer depuis la donnée d'origine, qui ne lui appartient pas. Aucune n'entre
donc dans un calcul du modèle et aucune ne porte un montant de capital. Ce qui est vérifiable
l'est néanmoins : les valeurs sont recopiées dans un script du dépôt, qui contrôle que le rapport
des sinistres aux primes redonne bien la série des ratios publiée, que la somme des classes de
taille redonne la charge annuelle sur les sept exercices, et que chaque multiplicateur cité se
retrouve à partir des niveaux qui l'encadrent. Ces contrôles portent sur la fidélité de la
recopie, non sur la collecte ni sur le périmètre de l'étude citée, qui relèvent de celle-ci.

\paragraph{Un marché qui s'élargit en baissant ses prix.} Entre $2024$ et $2025$, le nombre
d'entreprises assurées passe de $14\,124$ à $20\,996$, soit $+49\,\%$, quand les primes
souscrites reculent de $316{,}8$ à $305{,}9$~M\euro, soit $-3\,\%$ : l'élargissement du marché
s'est fait par la prime unitaire, qui perd $35\,\%$, et la franchise moyenne recule de
$149{,}9$ à $89$~k\euro. Dans le même temps, le nombre de sinistres indemnisés passe de $448$ à
$1\,251$.

\begin{figure}[htbp]
\centering
\figover{figures_externes/lucy_primes_sinistres_ratio.png}
\caption[Primes, sinistres et ratio du marché français de la cyberassurance]{Primes souscrites,
sinistres indemnisés et ratio sinistres sur primes du marché français de la cyberassurance,
$2019$ à $2025$, en millions d'euros. Reproduit du rapport LUCY 2026 de Nexialog Consulting
\citep{RapiorKaddouri2026}. L'étiquette du montant $2025$ des sinistres est recouverte par le
marqueur du ratio dans la source ; sa valeur est $83{,}2$~M\euro.}
\label{fig:lucy-sp}
\end{figure}

\noindent Le ratio sinistres sur primes remonte depuis trois exercices, de $12\,\%$ en $2023$ à
$17\,\%$ en $2024$ puis $27\,\%$ en $2025$, et l'écart double d'un exercice à l'autre, $5$ puis
$10$ points (figure~\ref{fig:lucy-sp}). Le marché a pourtant déjà connu un régime bien plus
dégradé, à $85\,\%$ en $2019$ et $169\,\%$ en $2020$ : le niveau de $2025$ est un point haut de
trois ans, non un point haut d'historique.

\paragraph{Deux régimes de risque sous une même trajectoire.} L'indicateur suivi par le marché
est le \emph{ratio sinistres sur primes}, c'est-à-dire le montant des sinistres payés dans
l'année rapporté aux primes encaissées la même année : il mesure la part de la prime qui part en
indemnisation. Ce ratio a monté deux exercices de suite, mais pour deux raisons opposées, et
seule sa décomposition le montre. La charge d'une année se met en effet sous la forme d'un
produit de deux termes, le \emph{nombre} de sinistres survenus et le \emph{sinistre moyen},
c'est-à-dire le coût moyen de l'un d'eux. Le tableau compare chaque terme à sa valeur de l'année
précédente, sous la forme d'un multiplicateur.

\begin{center}\small
\renewcommand{\arraystretch}{1.2}
\begin{tabular}{@{}l r r r l@{}}
\toprule
& \textbf{nombre} & \textbf{sinistre moyen} & \textbf{charge} & \textbf{régime} \\
\midrule
Marché, exercice $2024$ & $\times 0{,}73$ & $\times 1{,}96$ & $\times 1{,}43$ & sévérité \\
Marché, exercice $2025$ & $\times 2{,}79$ & $\times 0{,}55$ & $\times 1{,}53$ & fréquence \\
\midrule
\multicolumn{5}{@{}l}{\emph{Exercice $2025$, par bloc}} \\
grandes entreprises & $\times 1{,}24$ & $\times 0{,}82$ & $\times 1{,}02$ & absorbé \\
bloc intermédiaire & $\times 2{,}14$ & $\times 1{,}65$ & $\times 3{,}53$ & les deux \\
micro-entreprises & $\times 9{,}52$ & $\times 0{,}73$ & $\times 6{,}97$ & nombre seul \\
\bottomrule
\end{tabular}
\end{center}

\vspace{2pt}
\noindent En $2024$ le nombre de sinistres a baissé alors que leur
coût moyen a doublé : la charge a monté parce que les sinistres étaient plus lourds. En $2025$
c'est l'inverse, le coût moyen a baissé de moitié alors que le nombre a presque triplé : la
charge a monté parce que les sinistres étaient plus nombreux. Le ratio seul ne permet pas de
distinguer ces deux situations, qui n'appellent pourtant pas les mêmes mesures. Le bloc
intermédiaire est le seul segment où le nombre et le coût moyen augmentent en même temps, ce qui
explique sa charge multipliée par $3{,}53$. Le présent mémoire applique la même démarche à ses
quatre canaux de conformité au chapitre~\ref{chap:resultats} : un écart de capital se décompose
en ses causes au lieu de s'annoncer en niveau global.

\noindent La dégradation n'est pas uniforme entre segments (figure~\ref{fig:lucy-divergence}). Le
ratio des entreprises de taille intermédiaire passe de $13\,\%$ en $2024$ à $42\,\%$ en $2025$,
quand l'indice de leur taux de prime recule de son sommet de $257$ en $2023$ à $178$ en $2025$ :
le prix baisse sur le segment dont le risque se dégrade. Ce segment a déjà connu un ratio de
$261\,\%$ en $2021$.

\begin{figure}[htbp]
\centering
\figover{figures_externes/lucy_divergence_prix_risque.png}
\caption[Divergence entre prix et risque par segment]{Ratios sinistres sur primes par segment
(barres, axe de gauche) et indice du taux de prime des entreprises de taille intermédiaire
(courbe, axe de droite, base $100$ en $2020$). Pour chaque exercice, les trois barres sont, de
gauche à droite, les grandes entreprises, les entreprises de taille intermédiaire et les
entreprises moyennes. Reproduit du rapport LUCY 2026 de Nexialog Consulting
\citep{RapiorKaddouri2026}.}
\label{fig:lucy-divergence}
\end{figure}

\vspace{2pt}
\noindent Un point de vocabulaire accompagne ce tableau. La \emph{fréquence} désigne un nombre de
sinistres rapporté au nombre d'assurés, tandis que le \emph{nombre} de sinistres est un total. Ce
sont deux grandeurs distinctes : la fréquence vaut $1{,}88$ pour $2025$ quand le nombre vaut
$2{,}79$, la différence venant de l'\emph{exposition}, c'est-à-dire du volume d'assurés couverts,
qui a crû de $1{,}49$. L'égalité charge $=$ nombre $\times$ coût moyen n'est donc exacte qu'avec
le nombre. La distinction sépare ce qui relève d'une croissance du portefeuille de ce qui relève
d'une dégradation du risque à portefeuille constant.

\paragraph{Les sinistres les plus lourds sont trop rares en France pour être modélisés.} Cette
observation commande directement un choix du mémoire, car elle porte sur la région de la
distribution qui détermine le capital. Le capital de solvabilité se mesure en effet par un
\emph{quantile} à $99{,}5\,\%$, c'est-à-dire par le montant de perte annuelle qui n'est dépassé
qu'une année sur deux cents : il se joue donc entièrement dans les sinistres les plus lourds, et
non dans les sinistres courants. Or la donnée française y est presque vide. Deux exercices sur
sept ne portent aucun sinistre de la classe de taille la plus haute ; un exercice en porte
$135$~M\euro{} et l'exercice sous revue $19$~M\euro, pour un seul sinistre dépassant $10$~M\euro. Le
maximum de la série vaut $7{,}11$ fois le montant de l'exercice sous revue, ce qui donne la
mesure de l'irrégularité (figure~\ref{fig:lucy-taille}).

\begin{figure}[htbp]
\centering
\figover{figures_externes/lucy_taille_des_sinistres.png}
\caption[Montant indemnisé par classe de taille de sinistre]{Montant indemnisé par classe de
taille de sinistre sur le marché français, $2019$ à $2025$, en millions d'euros. Reproduit du
rapport LUCY 2026 de Nexialog Consulting \citep{RapiorKaddouri2026}. Les seuils de classe sont
ceux de la source.}
\label{fig:lucy-taille}
\end{figure}

\medskip
\noindent La conséquence est méthodologique. La méthode usuelle d'estimation d'une queue de
distribution, dite \emph{par dépassement de seuil}, consiste à ne retenir que les pertes
au-dessus d'un montant élevé et à ajuster une loi sur ces seules observations. Elle exige donc un
effectif suffisant au-dessus du seuil, et la donnée française n'en offre aucun : l'historique
national s'arrête avant le haut de la distribution, là précisément où se joue le quantile
recherché. Une estimation conduite sur cette base serait instable d'un exercice à l'autre et
sous-estimerait le risque, faute d'avoir observé les pertes qui le portent. C'est le motif,
repris au chapitre~\ref{chap:donnees}, de calibrer la sévérité sur une base internationale plutôt
que française. Cette rareté des événements extrêmes ne traduit pas une protection particulière du
marché français : les analyses de sinistres européennes relèvent un recul d'environ $20\,\%$ des
sinistres cyber en $2024$ qui laisse le niveau près d'un tiers au-dessus des exercices $2020$ à
$2022$ \citep{marsh2025}, et le marché britannique voisin est décrit comme plus développé et plus
exposé aux sinistres sévères \citep{dsit2025}.

\begin{attention}
\textbf{Ces montants ne se comparent pas à ceux du mémoire.} La charge citée ci-dessus est une
charge \emph{indemnisée} : c'est ce que les assureurs ont effectivement versé, donc après
déduction de la \emph{franchise} (la part du sinistre qui reste à la charge de l'assuré) et après
application de la \emph{capacité} du contrat (le montant maximal que l'assureur s'est engagé à
payer). Elle vaut $83{,}2$~M\euro{} en $2025$ contre $54{,}5$ l'exercice précédent, soit une
hausse de $53\,\%$. La sévérité que le présent mémoire modélise est d'une tout autre nature :
c'est la perte \emph{brute} subie par une entité financière, avant toute assurance. Les deux
grandeurs portent sur des périmètres différents et ne se comparent ni en niveau ni en quantile ;
les rapprocher conduirait à lire un écart de définition comme une contradiction.
\end{attention}

% SOURCES-SCRIPTS: 63

\section{Un seul fil, en cinq mouvements}
Le mémoire poursuit une seule question, \emph{combien coûte la non-conformité à DORA, et
jusqu'où ce chiffre peut-il être affirmé}, à travers cinq mouvements. Il construit d'abord le mécanisme, la cascade dirigée, et en démontre les propriétés
(chapitre~\ref{chap:cascade}). Il délimite ensuite ce que la
donnée permet d'en estimer : la frontière d'identifiabilité (chapitre~\ref{ch:identifiabilite}).
Il transforme cette limite en méthode, l'identification partielle, qui borne le
capital plutôt que de poser un paramètre inconnu (chapitre~\ref{chap:partielle}). Il fait alors
parler le modèle : la conformité devient un état qui évolue et pilote une trajectoire de capital
(chapitre~\ref{chap:multietats}). Il en tire enfin le surcoût, son attribution et l'ordre de
remédiation, et met en regard ce que les cadres existants savent en exprimer
(chapitre~\ref{chap:resultats}). Les deux derniers mouvements, la trajectoire et
les résultats, ne sont pas une thèse parallèle : ils sont l'\emph{objet} que les trois
premiers apprennent à lire avec la bonne prudence. Chaque niveau qu'ils affichent est
illustratif, chaque hiérarchie qu'ils proposent est ce qui reste vrai une fois la limite
d'identification prise en compte.

\section{Contribution et plan}
Le mémoire s'articule autour de trois apports, et il est aussi explicite sur ce qu'il
\emph{n'affirme pas} que sur ce qu'il établit.

\begin{enumerate}[leftmargin=1.4em,itemsep=3pt]
\item Un mécanisme de dépendance dirigée et à l'ordre entre les piliers, la cascade,
      dont trois propriétés sont démontrées (chapitre~\ref{chap:cascade},
      annexe~\ref{ann:cascade}) :
      la criticité dépend de l'\emph{ordre} de la chaîne et non du seul ensemble des piliers
      touchés, ce qu'aucune dépendance symétrique ne reproduit ; la normalisation de Leontief
      borne la propagation ; et une matrice et sa transposée, indistinguables pour la donnée,
      donnent un capital et une décision différents.
\item Une frontière d'identifiabilité, démontrée au
      chapitre~\ref{ch:identifiabilite}, puis muée en méthode au
      chapitre~\ref{chap:partielle} : la direction n'étant pas
      identifiable, le mémoire ne la pose pas, il borne le capital sur l'ensemble des matrices
      admissibles et chiffre en euros la valeur de la donnée manquante.
\item Une lecture décisionnelle : le surcoût $\Delta_{\text{DORA}}$, son attribution
      par pilier (Shapley, Euler) et l'ordre de remédiation, présentés en rapports et en
      bornes plutôt qu'en niveaux (chapitre~\ref{chap:resultats}).
\end{enumerate}

\begin{cle}
\textbf{Le résultat en bref.} Entre l'état conforme et l'état non conforme sur les quatre canaux
que la conformité déplace, le besoin de capital passe de $6\,049$ à $20\,188$~M\euro{} à
l'échelle du secteur, soit un facteur $3{,}34$ et un écart de $14\,139$~M\euro{}
(\S\ref{sec:chiffre-de-tete}). Ce montant est un besoin de capital au sens de l'ORSA et non un
SCR réglementaire : il sert à comparer deux états, non à remplacer la Formule Standard. Les
canaux interagissent, l'interaction portant $5\,001$~M\euro{}, soit $35\,\%$ de l'écart, et la
fréquence seule en produit davantage qu'aucun autre canal, à $10\,377$~M\euro{} contre
$11\,413$ pour les trois autres réunis. Les niveaux de propagation sont posés, mais la thèse ne
dépend que de leur ordre, le besoin de capital étant croissant en ce paramètre.

\textbf{Guide de lecture.} Un lecteur pressé trouve la thèse dans cette introduction, le chiffre
de tête et sa décomposition au chapitre~\ref{chap:resultats}, l'inventaire des hypothèses et
leur coût au chapitre~\ref{ch:robustesse}, et la validation hors échantillon au
\S\ref{soc:sec:backtest}. Les démonstrations et les pièces justificatives sont en annexe.
\end{cle}

\begin{attention}
\textbf{Ce que le mémoire ne prétend pas.} La matrice de contagion n'est jamais présentée
comme calibrée : elle est normative, posée à dire d'expert, et bornée par une analyse de
sensibilité. Le
niveau absolu du capital est illustratif, non une mesure : c'est l'effet \emph{relatif} de la
conformité et la hiérarchie des piliers qui sont défendus. Cette discipline est un
résultat du travail, et le fil directeur du document.
\end{attention}

% =====================================================================

% SOURCES-SCRIPTS: 29 43 68 92 66
